% answers_short.tex —— 2013 年·春季上海卷 · 简短答案速查
% 由 tools/make_short_answers.py 自 answers.tex 的 \gksolhead/\gkanswer 提取，
% 并与 2013/2013_数学真题_春季上海_答案要点.md 逐题比对后生成，勿手改（改源 MD 后重跑）。

\begin{gknotice}
\textbf{使用说明}：本节只给最终答案便于快速核对，完整推导见解析版。答案为编写者独立推导，\textbf{非官方标准答案}。
\end{gknotice}

\gkpart{填空题}
\begin{itemize}[leftmargin=2.6em]
  \item[\textbf{1.}] $(-2,+\infty)$。由真数 $x+2>0$ 得 $x>-2$。
  \item[\textbf{2.}] $3$。$2^x=8=2^3$，指数函数单调，故 $x=3$。
  \item[\textbf{3.}] $x=-2$。$y^2=8x$ 中 $2p=8$，准线为 $x=-\dfrac p2=-2$。
  \item[\textbf{4.}] $2\pi$。$\omega=1$，$T=\dfrac{2\pi}{|\omega|}=2\pi$。
  \item[\textbf{5.}] $-\dfrac34$。由 $1\cdot(k-6)-9k=0$ 得 $k=-\dfrac34$。
  \item[\textbf{6.}] $5$。$4\sin x+3\cos x=5\sin(x+\varphi)$，最大值为 $\sqrt{4^2+3^2}=5$。
  \item[\textbf{7.}] $\sqrt{13}$。$|2+3\mathrm i|=\sqrt{2^2+3^2}=\sqrt{13}$。
  \item[\textbf{8.}] $7$。由余弦定理 $b^2=5^2+8^2-2\cdot5\cdot8\cos60^\circ=49$；原卷印面“$b=8$”与所求 $b$ 矛盾，应为“$c=8$”。
  \item[\textbf{9.}] $\dfrac\pi3$。$\overrightarrow{A_1B}=(1,0,-1)$，$\overrightarrow{B_1C}=(0,1,-1)$，夹角余弦为 $\dfrac12$。
  \item[\textbf{10.}] $\dfrac45$。$1-\dfrac{\mathrm C_4^3+\mathrm C_6^3}{\mathrm C_{10}^3}=1-\dfrac{24}{120}=\dfrac45$。
  \item[\textbf{11.}] $\dfrac56n^2-\dfrac76n$。由 $6a_1+15d=23$、$9a_1+36d=57$ 得 $a_1=-\dfrac13$，$d=\dfrac53$。
  \item[\textbf{12.}] $4836$。$2000=2^4\cdot5^3$，约数和为 $(1+2+4+8+16)(1+5+25+125)=31\times156$。
\end{itemize}

\gkpart{选择题}
\begin{center}\begin{tabular}{ccccccccccccc}
\toprule
题号 & 13 & 14 & 15 & 16 & 17 & 18 & 19 & 20 & 21 & 22 & 23 & 24 \\
\midrule
答案 & B & B & D & A & D & A & C & B & C & A & D & C \\
\bottomrule
\end{tabular}\end{center}

\gkpart{解答题}
\begin{itemize}[leftmargin=2.6em]
  \item[\textbf{25.}] $18\sqrt3$。$BC=6\tan\dfrac\pi6=2\sqrt3$，$S_{\triangle ABC}=3\sqrt3$，$V=3\sqrt3\times6$。
  \item[\textbf{26.}] $500$ 平方米。设 $FP=x$，$BF=50-\dfrac54x$，面积 $=-\dfrac54(x-20)^2+500$。
  \item[\textbf{27.}] $\dfrac43$。$a_n=-2n+2$，$b_n=\left(\dfrac14\right)^{n-1}$，无穷等比级数和为 $\dfrac{1}{1-\frac14}$。
  \item[\textbf{28.}] (1) $\dfrac{3x^2}{4}+3y^2=1$（即 $\dfrac{x^2}{4/3}+\dfrac{y^2}{1/3}=1$）。由 $a=2b$ 与 $a^2-b^2=1$ 得 $b^2=\dfrac13$，$a^2=\dfrac43$；(2) $x+\sqrt7y-1=0$ 或 $x-\sqrt7y-1=0$。设 $y=k(x-1)$，由 $\overrightarrow{F_1P}\cdot\overrightarrow{F_1Q}=0$ 得 $k^2=\dfrac17$
  \item[\textbf{29.}] (1) $y^2=8-4x$。由 $\overrightarrow{AP}=-2\overrightarrow{FA}$ 得 $x_A=2-x$，$y_A=-y$，代入 $y_A^2=4x_A$；(2) $(0,0)$ 和 $\left(-\dfrac{15}{4},0\right)$。$Q(t,0)$ 的对称点为 $\left(-\dfrac{3t}5,\dfrac{4t}5\right)$，代入 $y^2=4x$ 得 $4t^2+15t=0$
  \item[\textbf{30.}] (1) $(0,4)$ 或 $(0,8)$。$\tan\theta_3=\dfrac{4t}{t^2+32}=\dfrac13$，解得 $t=4$ 或 $t=8$；(2) $n=4$ 时 $\theta_n$ 最大，最大值为 $\arctan\dfrac{\sqrt2}{4}$。$\tan\theta_n=\dfrac{1}{\frac{16\sqrt2}{2^n}+\frac{2^n}{8\sqrt2}}\leqslant\dfrac{\sqrt2}{4}$，等号在 $n=4$ 时成立
  \item[\textbf{31.}] (1) $y=x^3-3x$；对称中心为 $(1,-2)$。平移后为奇函数，对应 $a=1$，$b=-2$；(2) $(2,1)$。由定义域关于原点对称得 $a=2$，再由 $F(-x)+F(x)=0$ 得 $b=1$；(3) 假命题。反例 $f(x)=x$；修改为“函数 $y=f(x)$ 的图像关于直线 $x=a$ 成轴对称图像”的充要条件是“函数 $y=f(x+a)$ 是偶函数”
\end{itemize}
